\chapter{Derivations for Diffusion Models}
\label{app:diffusion-deriv}

Energy-based models, the DDPM variational bound, score matching in its two
equivalent forms, and the probability-flow ODE, supporting
Chapter~\ref{ch:background}.

\section{Energy-based models}\label{sec:ebm-def}

Modern generative modelling asks a single question: given samples of an
unknown distribution $p_0$, build a machine that produces new ones. The
last decade produced several answers: variational autoencoders
\citep{kingma2014auto}, adversarial networks
\citep{goodfellow2014generative}, normalizing flows
\citep{rezende2015variational}, diffusion models. It is easy to present
them as unrelated tricks. The energy-based view of
\citet{lecun2006tutorial} is the antidote: it is the direct continuation
of the Boltzmann-machine programme of
Section~\ref{sec:boltzmann-machines}, and, as this chapter unfolds,
diffusion models will appear as the member of the energy-based family
that finally solved its oldest problem.

An \emph{energy-based model} (EBM) over $x \in \mathcal{X}$ assigns to
each configuration a scalar energy $E_\theta(x)$ measuring how badly $x$
fits the model's idea of the data: low energy for plausible
configurations, high for implausible ones. In
\citeauthor{lecun2006tutorial}'s formulation the energy is the primitive
object: inference means finding low-energy configurations, and learning
means shaping the landscape, pushing the energy of observed data down and
the energy of everything else up. A probabilistic model is recovered
whenever one normalises with the Gibbs measure of
Chapter~\ref{ch:statmech},
\begin{equation}
  p_\theta(x) \;=\; \frac{e^{-E_\theta(x)}}{Z(\theta)},
  \qquad
  Z(\theta) = \int_{\mathcal{X}} e^{-E_\theta(x)}\, \dd x ,
  \label{eq:ebm}
\end{equation}
the Boltzmann distribution \eqref{eq:boltzmann} with $\beta$ absorbed
into the energy and the physics replaced by a parametric family. The
appeal of \eqref{eq:ebm} is that $E_\theta$ is unconstrained: any
function defines a valid unnormalised density. The price is the
partition function, and every basic operation hits it. \emph{Evaluating}
$p_\theta(x)$ requires $Z(\theta)$, intractable in general
(Section~\ref{sec:ising}). \emph{Sampling} requires the MCMC machinery of
Section~\ref{sec:mcmc}, whose mixing time can be exponential in the
presence of metastability. \emph{Learning} by maximum likelihood requires
the gradient
\begin{equation}
  -\nabla_\theta \log p_\theta(x)
  = \nabla_\theta E_\theta(x)
  - \E_{p_\theta}\!\big[ \nabla_\theta E_\theta \big],
  \label{eq:ebm-gradient}
\end{equation}
the ``push down on data, pull up on the model's own samples'' rule, the
continuous-variable form of the Boltzmann-machine rule
\eqref{eq:bm-learning}, whose second term again demands sampling from
$p_\theta$. \citet{lecun2006tutorial} frame the alternatives
systematically: either fight the negative phase with \emph{contrastive}
losses that pull up only on well-chosen offending configurations
(contrastive divergence \citep{hinton2002training} being the
MCMC-flavoured member of the family), or design architectures and
regularisers that bound the volume of low-energy space so that no pull-up
is needed. The lesson of the tutorial is that the partition function is
not an accounting detail; it is \emph{the} design constraint of the
field.

There are two classical escape routes, and this thesis walks both. The
first is the \emph{score}: since
$\nabla_x \log p_\theta(x) = -\nabla_x E_\theta(x)$, the gradient of the
log-density in $x$ is free of $Z$, because differentiating in $x$ kills
the constant. Sections~\ref{sec:score-matching}--\ref{sec:tweedie} build
the modern edifice standing on this observation. The second is
\emph{structure}: when the energy decomposes into local terms, marginals
and even $Z$ itself become computable by dynamic programming; that route
(Markov random fields, factor graphs, message passing) is developed in
Chapter~\ref{ch:graphical}.

\section{Denoising diffusion probabilistic models}\label{sec:ddpm}

The discrete-time formulation \citep{sohl2015deep, ho2020denoising}
discretises \eqref{eq:forward-marginal} into a Markov chain
$q(x_s | x_{s-1}) = \Nd(x_s; \sqrt{1 - \beta_s}\, x_{s-1}, \beta_s I)$
with a variance schedule $\beta_1, \dots, \beta_S$, and learns a reverse
chain $p_\theta(x_{s-1} | x_s)$. Training maximises a variational lower
bound on the data log-likelihood; the physicist will recognise the Gibbs
variational principle \eqref{eq:gibbs-variational} in modern dress.

\begin{toolbox}{The evidence lower bound of a diffusion model}
For any latent-variable model with joint $p_\theta(x_{0:S})$ and any
inference distribution $q(x_{1:S} | x_0)$, Jensen's inequality gives
\begin{align*}
  \log p_\theta(x_0)
  &= \log \int p_\theta(x_{0:S})\;\dd x_{1:S}
   = \log\, \E_{q}\!\left[
      \frac{p_\theta(x_{0:S})}{q(x_{1:S}|x_0)} \right] \\
  &\geq\; \E_{q}\!\left[
      \log \frac{p_\theta(x_{0:S})}{q(x_{1:S}|x_0)} \right]
  \;=:\; \mathrm{ELBO}.
\end{align*}
Insert the two Markov factorisations
$p_\theta(x_{0:S}) = p(x_S) \prod_s p_\theta(x_{s-1}|x_s)$ and
$q(x_{1:S}|x_0) = \prod_s q(x_s|x_{s-1})$, then regroup telescopically
using Bayes' rule
$q(x_s|x_{s-1}) = q(x_{s-1}|x_s, x_0)\, q(x_s|x_0) / q(x_{s-1}|x_0)$:
\begin{align*}
  -\mathrm{ELBO}
  \;=\;& \kl{q(x_S|x_0)}{p(x_S)}
  \;-\; \E_q \log p_\theta(x_0 | x_1) \\
  &+ \sum_{s=2}^{S}
    \E_q\, \kl{q(x_{s-1}|x_s, x_0)}{p_\theta(x_{s-1}|x_s)} ,
\end{align*}
one denoising step per summand. The decisive point: because the forward
chain is Gaussian and conditioned on \emph{both} endpoints, each target
$q(x_{s-1} | x_s, x_0)$ is an explicit Gaussian. Matching it with a
Gaussian $p_\theta$ reduces each KL term to a weighted regression of the
posterior mean; concretely, to predicting the noise $\varepsilon$ that
was added \citep{ho2020denoising}. The thermodynamic reading of the ELBO
gap as an entropy-production bound goes back to \citet{sohl2015deep}.
\hfill$\square$
\end{toolbox}

The practical upshot of \citet{ho2020denoising} is the loss
\begin{equation}
  \mathcal{L}(\theta)
  = \E_{s,\, a \sim p_0,\, \varepsilon}
    \big\| \varepsilon - \varepsilon_\theta\big(
      \underbrace{a\,e^{-t_s} + \sqrt{\Delta_{t_s}}\,\varepsilon}_{x_{t_s}},\; s
    \big) \big\|^2 ,
  \label{eq:ddpm-loss}
\end{equation}
a denoising regression, and the empirical discovery that this simplified
objective, with each noise level weighted equally and all the variational
bookkeeping dropped, trains photorealistic image models. Its connection
with the score is not incidental, as the next section shows: predicting
the noise \emph{is} estimating $\nabla \log p_t$ up to scale.

\section{Score matching}\label{sec:score-matching}

Suppose we model the score directly by a function $s_\theta(x)$ and wish
it to match $\nabla \log p(x)$ for data from $p$. The naive objective
$\E_p \| s_\theta - \nabla \log p \|^2$ seems to require knowing the very
score we lack. Two classical results remove the obstruction.

\begin{toolbox}{Implicit score matching (Hyv\"arinen)}
Expand the square and keep only $\theta$-dependent terms:
\begin{equation*}
  \E_p \norm{s_\theta - \nabla \log p}^2
  = \E_p \norm{s_\theta}^2
  - 2\, \E_p \inner{s_\theta}{\nabla \log p}
  + \text{const}.
\end{equation*}
The cross term integrates by parts. Componentwise, using
$p\,\partial_i \log p = \partial_i p$ and assuming
$p\, s_{\theta,i} \to 0$ at infinity:
\begin{equation*}
  \int p\; s_{\theta,i}\; \partial_i \log p \;\dd x
  = \int s_{\theta,i}\; \partial_i p \;\dd x
  = - \int p\; \partial_i s_{\theta,i} \;\dd x .
\end{equation*}
Therefore
\begin{equation*}
  \E_p \norm{s_\theta - \nabla \log p}^2
  = \E_p\Big[ \norm{s_\theta(x)}^2
    + 2\, \nabla \cdot s_\theta(x) \Big] + \text{const} :
\end{equation*}
an objective involving only the model and data samples
\citep{hyvarinen2005estimation}. Elegant, but the divergence term is
expensive for large networks, which is why diffusion models use the
denoising variant instead. \hfill$\square$
\end{toolbox}

\begin{toolbox}{Denoising score matching (Vincent)}
Fix $t$ and consider the noised marginal \eqref{eq:pt-convolution},
written as $p_t(x) = \int p_0(a)\, q_t(x|a)\,\dd a$ with Gaussian kernel
$q_t(x|a) = \Nd(x; a e^{-t}, \Deltat I)$. Then
\begin{equation*}
  \nabla_x \log p_t(x)
  = \frac{\int p_0(a)\, \nabla_x q_t(x|a)\, \dd a}{p_t(x)}
  = \int \underbrace{\frac{p_0(a)\, q_t(x|a)}{p_t(x)}}_{=\,p(a|x)}
    \nabla_x \log q_t(x|a)\, \dd a ,
\end{equation*}
so that
\begin{equation*}
  \nabla_x \log p_t(x)
  \;=\; \E\big[ \nabla_x \log q_t(x|a) \,\big|\, x \big] :
\end{equation*}
the score of the \emph{marginal} is the posterior average of the score of
the \emph{kernel}, and the latter is trivial,
$\nabla_x \log q_t(x|a) = -(x - a e^{-t})/\Deltat$. A standard
$L^2$-projection argument turns this into a regression: for any
$s_\theta$,
\begin{equation*}
  \E_{x \sim p_t} \norm{ s_\theta(x) - \nabla \log p_t(x) }^2
  = \E_{a, x} \norm{ s_\theta(x) - \nabla_x \log q_t(x|a) }^2
  + \text{const},
\end{equation*}
because conditional expectation is exactly the $L^2$ projection onto
functions of $x$ \citep{vincent2011connection}. Substituting the Gaussian
kernel score and $x = a e^{-t} + \sqrt{\Deltat}\varepsilon$ gives
$\nabla_x \log q_t = -\varepsilon / \sqrt{\Deltat}$: \emph{learning the
score of noisy data is learning to predict the noise}, and the DDPM loss
\eqref{eq:ddpm-loss} is denoising score matching in disguise.
\hfill$\square$
\end{toolbox}

The continuous-time synthesis \citep{song2019generative,
song2021scorebased} trains one network $s_\theta(x, t)$ against denoising
score matching at all noise levels and samples with the reverse SDE
\eqref{eq:reverse-sde}; an accessible exposition is \citet{song2021blog},
and \citet{karras2022elucidating} systematise the design space. This is
the resolution, promised in Section~\ref{sec:ebm-def}, of the
energy-based model's oldest problem: \citeauthor{song2019generative}
replace the energy by its gradient, the equilibrium sampler by a
finite-time reverse SDE, and the intractable likelihood by a regression
whose targets are exact.

\section{Flows: the deterministic side of the same coin}\label{sec:flows}

Anderson's theorem has a deterministic twin. As noted in
Section~\ref{sec:reversal}, the marginals $p_t$ of the forward SDE are
also transported by the noise-free \emph{probability-flow ODE} with
velocity $f - \tfrac{g^2}{2}\nabla \log p_t$ \citep{song2021scorebased}:
the same score, consumed by an ODE instead of an SDE, turns generation
into a smooth, invertible map from Gaussian noise to data. This places
diffusion models inside an older family. \emph{Normalizing flows}
\citep{rezende2015variational, papamakarios2021normalizing} construct an
invertible map $T_\theta$ and evaluate likelihoods through the
change-of-variables formula; their continuous-time version parametrises
the velocity field of an ODE directly \citep{chen2018neural}. What
diffusion brought to this family is a \emph{simulation-free training
signal}: \emph{flow matching} \citep{lipman2023flow} and stochastic
interpolants \citep{albergo2023building} show that one can regress the
velocity field of any prescribed interpolation between noise and data
(the OU path \eqref{eq:forward-marginal} being one choice among many)
with the same conditional trick as denoising score matching, never
integrating the ODE during training; a systematic introduction is
\citet{holderrieth2025flow}. The unifying statement is worth recording:
DDPMs, score SDEs, and flow-matching models all learn a vector field
whose ground truth is a \emph{conditional expectation under the noising
channel}; they differ in the path of marginals and in whether sampling is
stochastic or deterministic. Every exact result of this thesis therefore
speaks to the whole family at once: the object we compute in closed form,
the conditional expectation behind the field, is the training target of
all of them.

